%% toc: P3 33 33.1,33.2,33.3,33.4
\section{Further Material from Morley's \emph{Categoricity in Power}}
\label{sec:p03-w03-p03-033-morley-supplement}

This section completes the presentation of Morley's 1965 article begun in
\cref{sec:p03-n-a-006-01} with the parts of the article that sit beside its
main line of proof: the footnotes on terminology, on a syntactic reading
of the parameter systems, on open formulas and on Tarski's criterion, a
footnoted variant of the rank, the automorphism test
Morley uses to compute type spaces, a footnoted comparison between total
transcendence and the density of isolated types, the consequence and the
question about extending maps of an indiscernible set, the source's own
attributions, the footnotes on Jónsson's conditions, on algebraic
elements and on prime models over countable systems, and an example
asserted without proof in the footnote to Theorem~5.2. Every result below is stated for Morley's normalized theory
\(T=\Sigma^*\) of \cref{def:p03-morley-standing}: \(\Sigma\) is a complete
theory with an infinite model in an at most countable language, \(T\) is its
definitional relational expansion, \(\mathcal N(T)\) is the class of induced
subsystems of models of \(T\), and a monomorphism preserves and reflects
every relation and is injective. Types, the derivative
\(S^\alpha(A)\), \(\operatorname{Tr}^\alpha(A)\), rank and the degree
\(d(p)\) are those of \cref{def:p03-morley-type-space,def:p03-morley-global-rank}
and \cref{thm:p03-morley-degree}. An \emph{endomorphism} of a model \(A\) is
a monomorphism of \(A\) into \(A\), and an \emph{automorphism} is a
monomorphism of \(A\) onto \(A\); because \(T\) is Morleyized, both are
elementary on the original \(\Sigma\)-reduct.

The conclusions are these. Morley's concrete categories are categories in
the usual sense, parameter systems correspond to complete theories with
new constants, open formulas already give the whole formula algebra, and
satisfaction in a system agrees with provability from its theory for all
sentences exactly when the system is a model
(\cref{thm:p03-w03-p03-033-map-category,thm:p03-w03-p03-033-syntactic,thm:p03-w03-p03-033-open-formulas}).
A single ordinal rank defined by Morley's
one-word change of the derivative is monotone under extension and is
defined everywhere in a totally transcendental theory, but it does not
record the degree as a number (\cref{thm:p03-w03-p03-033-variant-rank,thm:p03-w03-p03-033-variant-field}).
A set of formulas lies in exactly one type when suitable automorphisms
exist, and a weakened automorphism condition is necessary as well as
sufficient (\cref{thm:p03-w03-p03-033-automorphism-test}). Density of
isolated types, the hypothesis under which prime models exist, is strictly
weaker than total transcendence (\cref{thm:p03-w03-p03-033-dlo-density}).
Totally transcendental theories have models in every uncountable size in
which every injection of a large set extends to an endomorphism
(\cref{thm:p03-w03-p03-033-endomorphism}), and Morley's question (5) has an
affirmative answer whenever a totally transcendental Skolem expansion is
available (\cref{thm:p03-w03-p03-033-question-five}). The class
\(\mathcal N(T)\) satisfies the closure conditions of Jónsson that the
article invokes, with amalgamation in its isomorphism form
(\cref{thm:p03-w03-p03-033-jonsson-conditions,thm:p03-w03-p03-033-jonsson-iv}).
Morley's algebraic points are exactly the types of Jónsson's algebraic
elements (\cref{thm:p03-w03-p03-033-jonsson-algebraic}). Over a countable
system a prime model exists under a hypothesis on that system alone, which
Morley's density hypothesis implies
(\cref{thm:p03-w03-p03-033-countable-prime}). Finally, a theory that is
not totally transcendental can lack countably saturated models in the
power \(\aleph_\omega\), with no assumption on the continuum
(\cref{thm:p03-w03-p03-033-coding-theory}). The opening problem of the
article, Łoś's question and its two framing examples, and the footnotes on
Jónsson's universal homogeneous systems and on Vaught's limit-power result
are treated in \cref{sec:p03-w03-p03-034-opening,sec:p03-w03-p03-034-universal-homogeneous}.
None of these results uses the later strongly minimal or pregeometry
analysis of \cref{sec:p03-n-a-006-02}.

\subsection{Footnotes on terminology, the syntactic reading, and references}
\label{sec:p03-w03-p03-033-early-footnotes}

Several of Morley's footnotes fix terminology, say how the semantic
set-up could be replaced by a syntactic one, or point to other works. They
are taken here in the order of the article.

\paragraph{Derived spaces (p.~515, n.~3).} Where the introduction
compares the new derivative with the familiar sequence of derived spaces
of a topological space, a footnote refers for that sequence to pages
126--134 of Kamke's \emph{Theory of Sets} (Dover, New York, 1950; entry
[7] of the article's reference list) \cite[p.~515, n.~3; p.~538]{Morley1965}.
The sequence meant is the Cantor--Bendixson sequence \(X^{(\alpha)}\) of
a space \(X\); \cref{thm:p03-w03-p03-034-cantor-bendixson} defines it and
proves that it is the one-space case of Morley's derivative. Kamke's pages
are not reconstructed here.

\paragraph{Categories of maps (p.~518, nn.~6 and 7).} Morley uses a
concrete notion of category. By a footnote, a \emph{mathematical object}
\(A\) is a set \(|A|\) together with some structure on it, and each
object occurring in the article is a relational system or a topological space
\cite[p.~518, n.~6]{Morley1965}.

\begin{hbdefinition}{def:p03-w03-p03-033-map-category}{Morley's category of maps}
Let \(\mathscr S\) be a class of mathematical objects. A \emph{category}
with object class \(\mathscr S\) is a class \(\mathscr C\) of triples,
called maps and written \(f:A\to B\), in which \(A,B\in\mathscr S\) and
\(f\) is a function of \(|A|\) into \(|B|\), such that (i) the identity
\(A\to A\) belongs to \(\mathscr C\) for each \(A\in\mathscr S\), and
(ii) whenever \((f:A\to B)\) and \((g:B\to C)\) belong to
\(\mathscr C\), so does \((gf:A\to C)\). In \(f:A\to B\), \(A\) is the
\emph{domain} and \(B\) the \emph{co-domain}
\cite[p.~518]{Morley1965}.
\end{hbdefinition}

The second footnote explains the choice. The more common definition
abstracts the way maps compose: a category is a class of elements with a
binary operation, defined only for certain pairs, subject to axioms. Morley
cares about particular categories rather than categories as such, and for
that purpose the concrete definition is easier to use
\cite[p.~518, n.~7]{Morley1965}. The next theorem records that nothing is
lost: a category in Morley's sense is a category in the abstract sense, and
the two categories of the article are categories in Morley's sense. Here
the abstract axioms are taken as follows: composition \(g\cdot f\) is
defined exactly when the co-domain of \(f\) is the domain of \(g\), it is
associative, and every object has an identity map that is a two-sided unit
for composition. By a footnote on p.~519, \(i_{AB}\) denotes the identity
map of \(A\) into \(B\) whenever \(A\subseteq B\)
\cite[p.~519, n.~11]{Morley1965}; it is the notation already used for the
restriction maps \(i_{AB}^*\) of \cref{def:p03-morley-type-space}.

\begin{hbtheorem}{thm:p03-w03-p03-033-map-category}{proved=annex:p03-w03-p03-033-map-category}{Morley's categories are categories}
\begin{enumerate}
\item If \(\mathscr C\) is a category in the sense of
\cref{def:p03-w03-p03-033-map-category}, then \(\mathscr C\), with
\((g:B\to C)\cdot(f:A\to B)=(gf:A\to C)\), satisfies the abstract axioms:
composition is defined exactly on pairs with matching co-domain and
domain, it is associative, and the identity triple of each object is a
two-sided unit.
\item The monomorphisms between members of \(\mathcal N(T)\) form a
category with object class \(\mathcal N(T)\).
\item Let \(\mathscr S(T)=\{S(A):A\in\mathcal N(T)\}\) and let
\(\mathscr C(T)\) be the class of triples \((f^*:S(B)\to S(A))\) for
monomorphisms \(f:A\to B\) in \(\mathcal N(T)\). Then
\(i_{AA}^*\) is the identity of \(S(A)\), \((gf)^*=f^*g^*\) for
monomorphisms \(f:A\to B\) and \(g:B\to C\), and \(\mathscr C(T)\) is a
category with object class \(\mathscr S(T)\) whose maps are continuous
and onto.
\end{enumerate}
\end{hbtheorem}

Morley describes \(\mathscr C(T)\) as dual to the category of
monomorphisms, and the formula \((gf)^*=f^*g^*\) of clause~3 is the
precise sense: the direction of every map is reversed
\cite[p.~519]{Morley1965}. Recording the co-domain in the triple matters.
A function alone does not determine its co-domain, so without it the
condition on composable pairs in clause~1 could not be read off a map.

\paragraph{The syntactic reading (p.~518, n.~8).} After defining
\(T(A)=T\cup\mathscr D(A)\), where the diagram \(\mathscr D(A)\) is the set
of open sentences of \(L(A)\) true in \((A,a)_{a\in|A|}\)
\cite[pp.~516--517]{Morley1965}, Morley notes in a footnote that the
whole of his Section~2 could have been written syntactically: in place of
the class \(\mathcal N(T)\) one would use the class of all complete
extensions of \(T\) in languages obtained from \(L\) by adding new
individual constants \cite[p.~518, n.~8]{Morley1965}. In the text he also
says that \(S(A)\) can be viewed as the set of ways of completing \(T(A)\)
in a language with one further individual constant
\cite[p.~518]{Morley1965}. The next theorem gives the dictionary that
makes both remarks exact. A complete theory is identified with the set of
its consequences. For a function \(f:|A|\to|B|\) and a formula \(\psi\)
of \(L(A)\), \(\tilde f\psi\) is the result of replacing each constant
\(\bar a\) by \(\overline{f(a)}\), as in \cref{def:p03-morley-type-space}.

\begin{hbtheorem}{thm:p03-w03-p03-033-syntactic}{proved=annex:p03-w03-p03-033-syntactic}{Parameter systems as complete theories with new constants}
\begin{enumerate}
\item For every \(A\in\mathcal N(T)\), \(T(A)\) is a consistent and
complete theory of \(L(A)\) that contains \(T\).
\item Let \(C\) be a set of individual constants not in \(L\) and let
\(T'\supseteq T\) be a complete theory of \(L\cup C\). Let \(M\) be a
model of \(T'\), let \(A\) be the subsystem of the \(L\)-reduct of \(M\)
on \(\{c^M:c\in C\}\), and let \(\sigma\) replace each \(c\in C\) by
\(\overline{c^M}\). Then \(A\in\mathcal N(T)\), and for every sentence
\(\varphi\) of \(L\cup C\), \(T'\vdash\varphi\) if and only if
\(T(A)\vdash\sigma\varphi\). If \(M'\) is another model of \(T'\) with
corresponding system \(A'\), the assignment \(c^M\mapsto c^{M'}\) is a
well-defined isomorphism of \(A\) onto \(A'\).
\item For \(A,B\in\mathcal N(T)\) and a function \(f:|A|\to|B|\), the
following are equivalent: \(f\) is a monomorphism; \(T(B)\vdash\tilde
f\psi\) for every \(\psi\in\mathscr D(A)\); \(T(B)\vdash\tilde f\psi\)
for every sentence \(\psi\) with \(T(A)\vdash\psi\).
\item Let \(c\) be an individual constant not in \(L(A)\), and for
\(p\in S(A)\) let \(p(c)\) be the set of sentences obtained from the
formulas of \(p\) by substituting \(c\) for \(v_0\). Then
\(p\mapsto T(A)\cup p(c)\) is a one-one map of \(S(A)\) onto the set of
complete theories of \(L(A)\cup\{c\}\) that contain \(T(A)\).
\end{enumerate}
\end{hbtheorem}

Under this dictionary a member of \(\mathcal N(T)\) is recovered, up to
isomorphism, from a complete extension of \(T\) by new constants, a
monomorphism becomes an inclusion of theories after renaming constants,
and a point of \(S(A)\) becomes a complete extension of \(T(A)\) by one
constant. The proof of clause~1 uses the normalization
\(T=\Sigma^*\) through Lemma~1.1(e); the other clauses use only clause~1
and the definitions.

\paragraph{Open formulas (p.~518, n.~9).} A further footnote makes two
observations about \(F(A)\). By Lemma~1.1(e), the Boolean algebra would be
the same if \(F(A)\) contained open formulas only. For open formulas,
satisfaction in \(A\) (written \(\vdash_A\) in the source) agrees with
provability from \(T(A)\), but for arbitrary formulas the two agree only
when \(A\) is a model of \(T\) \cite[p.~518, n.~9]{Morley1965}. The
second symbol of the footnote is printed as \(\vdash_{T(A)}\) with an
accent over the \(A\); it is read here as provability from \(T(A)\). Below,
\(A\entails\psi\) means that the sentence \(\psi\) of \(L(A)\) holds in
\((A,a)_{a\in|A|}\), which is Morley's \(\vdash_A\psi\)
\cite[p.~516]{Morley1965}. Lemma~1.1(e) is the clause of
\cref{thm:p03-morley-normalization} saying that every formula is
equivalent modulo \(T\) to an atomic relation formula.

\begin{hbtheorem}{thm:p03-w03-p03-033-open-formulas}{proved=annex:p03-w03-p03-033-open-formulas}{Open formulas suffice, and satisfaction agrees with provability only on them}
Let \(A\in\mathcal N(T)\).
\begin{enumerate}
\item Every \(\psi\in F(A)\) is equivalent in \(T(A)\) to an open formula
\[
 R(v_0,\bar a_1,\dots,\bar a_n),
\]
where \(R\) is a relation symbol of
\(L\) and \(a_1,\dots,a_n\in|A|\). Hence the open formulas of \(F(A)\),
with formulas identified when they are equivalent in \(T(A)\), form a
Boolean algebra that the inclusion maps isomorphically onto \(F(A)\), and
the two Stone spaces coincide.
\item For every open sentence \(\psi\) of \(L(A)\),
\(A\entails\psi\) if and only if \(T(A)\vdash\psi\).
\item If \(A\) is a model of \(T\), then \(A\entails\psi\) if and only if
\(T(A)\vdash\psi\), for every sentence \(\psi\) of \(L(A)\).
\item If \(A\) is not a model of \(T\), some sentence \(\psi\) of \(L\)
satisfies \(T(A)\vdash\psi\) and \(A\not\entails\psi\).
\end{enumerate}
\end{hbtheorem}

\begin{hbexample}{def:p03-w03-p03-033-open-formulas}{Two points of a dense order}
Let \(\Sigma=\mathsf{DLO}\), \(T=\Sigma^*\), and let \(A\) be the
subsystem of a model of \(T\) on two elements \(a<b\). The sentence
\(\psi=\exists v_0(\bar a<v_0\wedge v_0<\bar b)\) is provable from
\(T(A)\), because \(\bar a<\bar b\) is in \(\mathscr D(A)\) and \(T\)
says the order is dense, but \(A\not\entails\psi\), since \(A\) has no
element strictly between \(a\) and \(b\). The relation symbol \(R\) that
\(T\) attaches to the formula \(\exists v_2(v_0<v_2\wedge v_2<v_1)\)
holds of \((a,b)\) in \(A\), so the open sentence \(R(\bar a,\bar b)\)
is true in \(A\) and provable from \(T(A)\), as clause~2 requires. The
quantified sentence and the open sentence are equivalent in \(T\) but not
in \(A\), which is not a model of \(T\). These facts are verified at the
end of the proof of \cref{thm:p03-w03-p03-033-open-formulas}.
\end{hbexample}

\paragraph{Boolean algebras of formulas (p.~518, n.~10).} Morley's
footnote to the Boolean algebra \(F(A)\) says that several authors had
observed the close connection between properties of the Boolean algebras
of formulas of \(L\) and model-theoretic properties of \(T\), and names
especially Ryll-Nardzewski's 1959 paper on theories categorical in power
at most \(\aleph_0\) and Vaught's 1961 \emph{Denumerable models of
complete theories} (entries [13] and [18]) \cite[p.~518, n.~10;
p.~538]{Morley1965}. The handbook states Ryll-Nardzewski's theorem in
\cref{sec:qe-categoricity-02} and presents its original finite-type
criterion in \cref{sec:p03-n-a-001-01} \cite{RyllNardzewski1959}; it presents
Vaught's paper in \cref{sec:p03-n-a-003-01} \cite{Vaught1961}.

\paragraph{Algebraically closed fields (p.~523, n.~15).} For Example~I,
the theory of algebraically closed fields of characteristic \(0\)
(\cref{def:p03-morley-field-data,thm:p03-morley-field-example}), Morley
refers the reader to Abraham Robinson's \emph{Complete Theories} for a
fuller discussion \cite[p.~523, n.~15]{Morley1965}
\cite{Robinson1956}. No result below depends on that book.

\paragraph{Tarski's criterion (p.~529, n.~19).} In the proof of his
Lemma~4.1, that a system is a model exactly when its realized points are
dense in \(S(A)\) (the first clause of \cref{thm:p03-morley-prime-models}),
Morley invokes a theorem of Tarski, identified in a footnote as
Theorem~1.10 of his reference [15], Tarski and Vaught's
\emph{Arithmetical extensions of relational systems} (Compositio
Mathematica 13, 1957) \cite[p.~529, n.~19; p.~538]{Morley1965}. The
criterion as Morley uses it reads as follows; an elementary monomorphism is
a map \(f\) of \(X\subseteq|A|\) into \(|B|\) with
\((A,x)_{x\in X}\equiv(B,f(x))_{x\in X}\) \cite[p.~517]{Morley1965}.

\begin{hbcorollary}{thm:p03-w03-p03-033-tarski-criterion}{proved=annex:p03-w03-p03-033-tarski-criterion}{Tarski's criterion in Morley's form}
Let \(B\) be a model of \(T\) and \(A\subseteq B\) with
\(A\in\mathcal N(T)\) and \(|A|\) nonempty. Then \(i_{AB}\) is an elementary monomorphism if
and only if every formula \(\psi\in F(A)\) that some \(b\in|B|\)
satisfies in \(B\) is satisfied in \(B\) by some \(a\in|A|\). In that case
\(A\) is a model of \(T\).
\end{hbcorollary}

The corollary is the Tarski--Vaught test (\cref{thm:tarski-vaught-fol})
written in Morley's notation; it is not taken from the Tarski--Vaught
paper, which is not reconstructed here. Satisfaction is in \(B\) on both
sides. In \cref{def:p03-w03-p03-033-open-formulas} the formula
\(\bar a<v_0\wedge v_0<\bar b\) is satisfied in a model \(B\supseteq A\)
by an element between \(a\) and \(b\) and by no element of \(A\), so
\(i_{AB}\) is not elementary, in agreement with clause~4 of
\cref{thm:p03-w03-p03-033-open-formulas}.

\subsection{A single rank in place of rank and degree}
\label{sec:p03-w03-p03-033-variant-rank}

In a footnote to his Theorem~2.5 on degree, Morley remarks that rank and degree can
be merged into one new rank by a slight change in the definition of the
derivative: in clause (2) of his Definition~2.2 the requirement that the
extension fibres form a set of isolated points is replaced by the
requirement that they form a single point
\cite[p.~521, n.~13]{Morley1965}. He gives no proof and does not say which
properties the merged rank has. The reading adopted here takes ``a single
point'' to mean ``at most one point''. The worked example below shows that
the reading with ``exactly one point'' leaves a type of an algebraically
closed field without any rank, although that theory is totally
transcendental.

\begin{hbdefinition}{def:p03-w03-p03-033-variant-rank}{The single-point derivative}
Simultaneously for all \(A\in\mathcal N(T)\) and by recursion on the
ordinal \(\beta\), put
\[
 S_\natural^\beta(A)=S(A)\setminus\bigcup_{\gamma<\beta}
                  \operatorname{Tr}_\natural^\gamma(A),
\]
and let \(p\in\operatorname{Tr}_\natural^\beta(A)\) when
\(p\in S_\natural^\beta(A)\) and, for every monomorphism
\(f:A\to B\) with \(B\in\mathcal N(T)\), the set
\((f^*)^{-1}(p)\cap S_\natural^\beta(B)\) has at most one point.
The \emph{single-point rank} \(\rho_\natural(p)\) is the unique
\(\beta\) with \(p\in\operatorname{Tr}_\natural^\beta(A)\); a point that
lies in every \(S_\natural^\beta(A)\) has no single-point rank.
The only change from \cref{def:p03-morley-global-rank} is the fibre
condition: Morley's derivative asks that every fibre point be isolated in
\(S^\beta(B)\), whereas here the fibre must have at most one point and no
isolation is asked.
\end{hbdefinition}

Uniqueness of \(\beta\) holds because \(\operatorname{Tr}_\natural^\beta(A)\)
is contained in \(S_\natural^\beta(A)\), which is disjoint from every
earlier \(\operatorname{Tr}_\natural^\gamma(A)\).

\begin{hbtheorem}{thm:p03-w03-p03-033-variant-rank}{proved=annex:p03-w03-p03-033-variant-rank}{What the single-point rank combines}
Let \(f:A\to B\) be a monomorphism in \(\mathcal N(T)\) and let
\(p\in S(A)\).
\begin{enumerate}
\item If \(p\) has single-point rank \(\beta\), every point of
\((f^*)^{-1}(p)\) has a single-point rank at most \(\beta\), and at most
one point of \((f^*)^{-1}(p)\) has single-point rank exactly \(\beta\).
\item Restriction does not lower membership: if
\(q\in S_\natural^\beta(B)\), then \(f^*(q)\in S_\natural^\beta(A)\).
\item If \(T\) is totally transcendental, every point of every
\(S(A)\) has a single-point rank.
\end{enumerate}
\end{hbtheorem}

Clause~1 is the combined role that rank and degree play together in
Morley's arguments: by \cref{thm:p03-morley-degree}, an extension of a
point has no larger rank, and among extensions of equal rank the degree
cannot increase, with equality for at most one extension. The single
ordinal \(\rho_\natural\) gives the same two conclusions at once. The proof
of clause~3 uses the pair (rank, degree) of \cref{thm:p03-morley-degree}
in lexicographic order.

\begin{hbexample}{def:p03-w03-p03-033-variant-field}{Two algebraic types in characteristic zero}
Let \(\Sigma\) be the complete theory of algebraically closed fields of
characteristic \(0\), as in \cref{def:p03-morley-field-data}, and let
\(A=\{0\}\), so the field generated by \(A\) is \(\mathbb Q\). Let
\(p_2\in S(A)\) be the type containing \(v_0^2=2\) and \(p_3\in S(A)\)
the type containing \(v_0^3=2\). The polynomials \(v^2-2\) and
\(v^3-2\) are irreducible over \(\mathbb Q\) and have \(2\) and \(3\)
distinct roots, so by \cref{thm:p03-morley-field-example} both types have
rank \(0\), with degrees \(d(p_2)=2\) and \(d(p_3)=3\).
\end{hbexample}

\begin{hbtheorem}{thm:p03-w03-p03-033-variant-field}{proved=annex:p03-w03-p03-033-variant-field}{Both types have single-point rank one}
In \cref{def:p03-w03-p03-033-variant-field},
\(\rho_\natural(p_2)=\rho_\natural(p_3)=1\). Hence the single-point rank
does not determine the degree. If the fibre condition of
\cref{def:p03-w03-p03-033-variant-rank} is read as ``exactly one point'',
\(p_3\) lies in every stage of the modified derivative and so has no rank,
although \(\mathrm{ACF}_0\) is totally transcendental.
\end{hbtheorem}

\begin{hbwarning}{def:p03-w03-p03-033-variant-limits}{What the single-point rank does not do}
Morley's derivative maps each stage onto the same stage under restriction
(\cref{thm:p03-morley-derivative}); the single-point derivative does not.
In the example, over a system containing all three roots of \(v^3-2\),
every extension of \(p_3\) has single-point rank \(0\), although
\(\rho_\natural(p_3)=1\). The rank is therefore not a numerical encoding of
the pair (rank, degree): degrees \(2\) and \(3\) receive the same value.
\Cref{thm:p03-w03-p03-033-variant-rank} records what the merged rank
does provide. Its comparison with Morley rank in general, and the converse
of clause~3, are not established here; the source states neither.
\end{hbwarning}

\subsection{An automorphism test for unique types}
\label{sec:p03-w03-p03-033-automorphism-test}

Before his three examples Morley needs to know when a consistent set of
formulas determines a type. He gives a sufficient condition in terms of
automorphisms of models, proves it, and adds in a footnote that a
weakened form of the condition is also necessary, referring to the
Morley--Vaught paper on homogeneous universal models
\cite[p.~523 and n.~14]{Morley1965}. Here \(F(A)\) is the formula algebra
of \cref{def:p03-morley-type-space}, a set \(Q\subseteq F(A)\) is
consistent when some point of \(S(A)\) contains it, and an element
\(b\) of a system containing \(A\) satisfies \(Q\) when every formula of
\(Q\) holds of \(b\).

\begin{hbtheorem}{thm:p03-w03-p03-033-automorphism-test}{proved=annex:p03-w03-p03-033-automorphism-test}{Automorphisms decide uniqueness of a type}
Let \(A\in\mathcal N(T)\) and let \(Q\subseteq F(A)\) be consistent.
\begin{enumerate}
\item (Morley's condition.) Suppose that whenever \(B\) is a model of
\(T\) with \(A\subseteq B\) and \(b,b'\in B\) satisfy \(Q\), some
automorphism of \(B\) fixes every element of \(A\) and carries \(b\) to
\(b'\). Then \(Q\) is contained in exactly one point of \(S(A)\).
\item (The weakened condition.) Say that \(Q\) has the weakened property
when, for every model \(B\supseteq A\) of \(T\) and all
\(b,b'\in B\) satisfying \(Q\), there are a model \(C\) of \(T\) with
\(B\subseteq C\) and an automorphism of \(C\) that fixes every element
of \(A\) and carries \(b\) to \(b'\). Then \(Q\) is contained in exactly
one point of \(S(A)\) if and only if \(Q\) has the weakened property.
\end{enumerate}
\end{hbtheorem}

Morley's condition implies the weakened one (take \(C=B\)), so clause~1
follows from the sufficiency half of clause~2. Morley applies the
condition to his first example: an irreducible polynomial equation over the
field generated by \(A\), and the set of all non-vanishing conditions,
each lie in exactly one type. The proof of
\cref{thm:p03-morley-field-example} uses automorphisms of a common
algebraically closed extension, which is the weakened form.

\subsection{Dense isolated types without total transcendence}
\label{sec:p03-w03-p03-033-density}

Prime models over arbitrary subsystems exist as soon as the isolated
points of every \(S(A)\) are dense, and total transcendence implies this
density (\cref{thm:p03-morley-prime-models}). In a footnote to his chain
theorem, Theorem~4.5, Morley records three things \cite[p.~532, n.~21]{Morley1965}:
an example shows that total transcendence is a strictly stronger
assumption than density of the isolated points in every \(S(A)\); the
prime-model theorem, his Theorem~4.3 (\cref{thm:p03-morley-prime-models}), was proved under the weaker density assumption; and
he could not prove Theorem~4.5 on continuous chains of prime models
(\cref{thm:p03-morley-chain-primes}) under it. The footnote does not
name the example. Dense linear order without endpoints, which Morley
uses elsewhere to show that countable categoricity does not give total
transcendence \cite[p.~528]{Morley1965}, separates the two assumptions.

\begin{hbtheorem}{thm:p03-w03-p03-033-dlo-density}{proved=annex:p03-w03-p03-033-dlo-density}{Density of isolated types is strictly weaker}
Let \(\Sigma=\mathsf{DLO}\), the theory of dense linear orders without
endpoints, and \(T=\Sigma^*\). For every \(A\in\mathcal N(T)\), including
the empty system, the isolated points of \(S(A)\) are dense. The theory
\(T\) is not totally transcendental. Consequently every subsystem of a
model of \(\mathsf{DLO}\) has a prime model over it, while the hypothesis
of \cref{thm:p03-morley-chain-primes} fails.
\end{hbtheorem}

The theorem shows only that the hypothesis of the chain theorem is not
met. Whether continuous chains of prime models exist for every theory
with dense isolated types is the question Morley leaves open in that
footnote; nothing here decides it.

\subsection{Extending maps of an indiscernible set}
\label{sec:p03-w03-p03-033-question-five}

After the four questions about categoricity Morley returns to countable
languages and derives a consequence of his Theorem~4.3 on prime models and his
Theorem~4.6 on the extraction of indiscernible sets
(\cref{thm:p03-morley-prime-models,thm:p03-morley-indiscernible-extraction})
\cite[p.~537]{Morley1965}.

\begin{hbtheorem}{thm:p03-w03-p03-033-endomorphism}{proved=annex:p03-w03-p03-033-endomorphism}{Every injection of a large set extends to an endomorphism}
Let \(T\) be totally transcendental and \(\kappa>\aleph_0\). There are a
model \(A\) of \(T\) and a set \(X\subseteq|A|\) with
\(\kappa(X)=\kappa(A)=\kappa\) such that every one-one map of \(X\) into
itself extends to an endomorphism of \(A\).
\end{hbtheorem}

Morley then asks question~(5): if \(T\) is totally transcendental and
\(\kappa\geq\aleph_0\), is there always a model \(A\) with a set
\(X\subseteq|A|\), \(\kappa(A)=\kappa(X)=\kappa\), such that every one-one
map of \(X\) onto itself extends to an automorphism of \(A\)? He notes that
the answer would follow from his Theorem~3.5 (extension of order
automorphisms to term hulls) and Theorem~3.9 (exclusion of connected
antisymmetric relations,
\cref{thm:p03-morley-no-order}) if a totally transcendental \(T\) always had a totally transcendental
Skolem expansion \cite[p.~537]{Morley1965}. The next theorem proves that
conditional. It does not assert that such an expansion exists.

\begin{hbtheorem}{thm:p03-w03-p03-033-question-five}{proved=annex:p03-w03-p03-033-question-five}{Question (5) under a totally transcendental Skolem expansion}
Let \(T^\#\) be a Skolem expansion of \(T\) in the sense of
\cref{def:p03-morley-skolem-hull}, in a countable language
\(L^\#\), and let \(T'\supseteq T^\#\) be a complete \(L^\#\)-theory whose
Morleyization is totally transcendental. Then for every
\(\kappa\geq\aleph_0\) there are a model \(A\) of \(T\) and a set
\(X\subseteq|A|\) with \(\kappa(A)=\kappa(X)=\kappa\) such that every
one-one map of \(X\) onto itself extends to an automorphism of \(A\).
\end{hbtheorem}

Morley also reports that Theorem~2 of Ehrenfeucht's 1957 paper
\emph{On theories categorical in power} (his reference [1]) asserts the
affirmative answer to (5) for theories categorical in a power
\(2^\kappa\), and that Vaught found a fallacy in the proof given there
\cite[p.~537]{Morley1965}. The theorem above makes no use of that claim.
Morley's footnote to the list of questions adds that Jack Silver had
recently answered (5) affirmatively \cite[p.~537, n.~25]{Morley1965}; the
present section does not reconstruct Silver's argument.

\subsection{Attributions recorded in the source}
\label{sec:p03-w03-p03-033-attributions}

Morley's footnotes assign several results and questions to others. The
numbers in brackets are entries of the article's own reference list.
\begin{itemize}
\item For the parameter-free case \(X=\varnothing\) of the thin-model
theorem (\cref{thm:p03-morley-thin-models}), Morley credits Ehrenfeucht
[2]. Ehrenfeucht obtained more: if two elements are called equivalent
when an automorphism of the model carries one to the other, a model of
power \(\kappa\) still exists with only countably many equivalence
classes. Morley notes that the proof resembles his own but needs a more
elaborate linear ordering of the generators
\cite[p.~528, n.~16]{Morley1965}.
\item The key step in deriving total transcendence from categoricity,
that \(S(C)\) is countable for every countable \(C\in\mathcal N(T)\),
had been established by Vaught, in the case of categoricity in a power
\(\kappa=\kappa^{\aleph_0}\); the footnote points to [10], the
Morley--Vaught paper \cite[p.~528, n.~17]{Morley1965}.
\item The exclusion of connected antisymmetric relations
(\cref{thm:p03-morley-no-order}) had been obtained by Ehrenfeucht [1]
for theories categorical in a power \(2^\kappa\). Dana Scott, in
unpublished work and by a different and simpler proof, extended it to
theories categorical in a power \(\kappa^{\aleph_0}\)
\cite[p.~528, n.~18]{Morley1965}.
\item Problems (1)--(4) in the closing list (\cref{sec:p03-morley-closing-questions})
are not Morley's own: he says they appear to have been considered by
several people \cite[p.~537, n.~25]{Morley1965}.
\end{itemize}

These attributions report Morley's account in 1965. The section does not
check them against the works cited, which are not reconstructed here.

\subsection{Jónsson's conditions for the class of parameter systems}
\label{sec:p03-w03-p03-033-jonsson-conditions}

Morley's normalization theorem collects three closure properties of
\(\mathcal N(T)\): unions of chains, joint embedding, and amalgamation over
a common subsystem (\cref{thm:p03-morley-normalization}). A footnote
explains where they come from \cite[p.~517, n.~5]{Morley1965}. Jónsson
had studied classes of relational systems subject to conditions numbered
I to VI; in order to apply his results to an arbitrary complete theory,
the Morley--Vaught paper introduced the theory \(\Sigma^*\) and showed
that \(\mathcal N(\Sigma^*)\) satisfies those conditions; and Morley's
clauses 1.2(b), (c) and (d) say respectively that \(\mathcal N(\Sigma^*)\)
satisfies Jónsson's conditions V, III and IV. The Jónsson paper cited by
Morley names the conditions by reference to Jónsson's earlier paper
\cite[pp.~137--138]{Jonsson1960}, where they are listed
\cite[p.~195]{Jonsson1956}; the later paper replaces IV by a weaker
condition IV\('\) and I by a stronger condition I\('\)
\cite[pp.~137--138]{Jonsson1960}. Jónsson's systems have a nonempty
universe and a family of relations of finite rank, finite in the earlier
paper and at most denumerable in the later one. Jónsson says only that
isomorphism is defined in the obvious manner \cite[p.~194]{Jonsson1956};
here an isomorphism of one system into another is read as a one-one map
that preserves and reflects each relation, which for \(\mathcal N(T)\) is
Morley's monomorphism. For
systems \(A,B\), \(A\prec B\) means that \(A\) is a subsystem of \(B\),
and \(^*A\) is the power of \(A\).

\begin{hbdefinition}{def:p03-w03-p03-033-jonsson-conditions}{Jónsson's conditions on a class of systems}
Let \(K\) be a class of relational systems of one similarity type and let
\(\alpha\) be an ordinal. The conditions, in Jónsson's numbering
\cite[p.~195]{Jonsson1956} and \cite[pp.~137--138]{Jonsson1960}, are:
\begin{description}
\item[I] there are \(A,B\in K\) that are not isomorphic;
\item[I\('\)] for each ordinal \(\xi\) there is \(A\in K\) with \(^*A\geq\aleph_\xi\);
\item[II] if \(A\in K\) and \(A\) is isomorphic to \(B\), then \(B\in K\);
\item[III] for all \(A,B\in K\) there is \(C\in K\) such that \(A\) and
\(B\) are isomorphic to subsystems of \(C\);
\item[IV] for all \(A,B\in K\), if \(A\cap B\in K\), \(A\cap B\prec A\)
and \(A\cap B\prec B\), there is \(C\in K\) with \(A\prec C\) and
\(B\prec C\);
\item[IV\('\)] if \(f_0\) and \(f_1\) are isomorphisms of \(A\in K\)
into \(B_0\in K\) and \(B_1\in K\), there are \(C\in K\) and
isomorphisms \(g_0\) of \(B_0\) and \(g_1\) of \(B_1\) into \(C\) with
\(g_0f_0(x)=g_1f_1(x)\) for every element \(x\) of \(A\);
\item[V] if \(\lambda\) is a positive ordinal, \(A_\xi\in K\) for every
\(\xi<\lambda\), and \(A_\xi\prec A_\eta\) whenever \(\xi<\eta<\lambda\),
then \(\bigcup_{\xi<\lambda}A_\xi\in K\);
\item[VI\(_\alpha\)] if \(A\in K\), \(B\prec A\) and \(^*B<\aleph_\alpha\),
there is \(C\in K\) with \(B\prec C\prec A\) and \(^*C<\aleph_\alpha\).
\end{description}
\end{hbdefinition}

\begin{hbtheorem}{thm:p03-w03-p03-033-jonsson-conditions}{proved=annex:p03-w03-p03-033-jonsson-conditions}{The parameter systems satisfy Jónsson's conditions}
The class \(K=\mathcal N(T)\) satisfies I, I\('\), II, III, IV\('\), V and
VI\(_\alpha\) for every ordinal \(\alpha\). Clause V is the union clause
of \cref{thm:p03-morley-normalization}, III is its joint embedding clause,
and IV\('\) is its amalgamation clause. Moreover \(\mathcal N(T)\) is
closed under subsystems and under unions of nonempty directed families.
\end{hbtheorem}

Morley's Lemma~1.2(d) has the isomorphism form IV\('\), which is the form
used in the Jónsson paper Morley cites. The original subsystem form IV is
a different condition, and it fails for some complete theories.

\begin{hbexample}{def:p03-w03-p03-033-jonsson-iv}{Successor in the integers}
Let \(\Sigma\) be the complete theory of \((\mathbb Z,<)\) and
\(T=\Sigma^*\). Let \(A\) and \(B\) be two-element systems with
\(|A|=\{0,a\}\) and \(|B|=\{0,b\}\), where \(a\neq b\) and neither is
\(0\), each made isomorphic to the subsystem of the expansion of
\((\mathbb Z,<)\) on \(\{0,1\}\) by the map that fixes \(0\) and sends
\(a\), respectively \(b\), to \(1\).
\end{hbexample}

\begin{hbtheorem}{thm:p03-w03-p03-033-jonsson-iv}{proved=annex:p03-w03-p03-033-jonsson-iv}{Condition IV fails for the integers}
In \cref{def:p03-w03-p03-033-jonsson-iv}, \(A,B\in\mathcal N(T)\),
\(A\cap B\) is the one-element system on \(\{0\}\), which belongs to
\(\mathcal N(T)\) and is a subsystem of both, and no \(C\in\mathcal N(T)\)
has both \(A\) and \(B\) as subsystems. Hence \(\mathcal N(T)\) does not
satisfy condition IV, although it satisfies IV\('\).
\end{hbtheorem}

Jónsson proved that, under conditions II, V and VI\(_\alpha\), any two
systems of power \(\aleph_\alpha\) that are universal and homogeneous for
\(K\) are isomorphic \cite[p.~139, Theorem~B]{Jonsson1960}. By Morley's
account, the Morley--Vaught paper obtained uniqueness of saturated models
from a more general result of this paper \cite[p.~514, n.~2]{Morley1965}.
The uniqueness theorem used in this chapter is proved directly in
\cref{thm:p03-morley-saturation-existence}, and the identification of
saturated with universal homogeneous systems is
\cref{thm:p03-w03-p03-034-universal-homogeneous}; neither depends on
Jónsson's theorem.

\subsection{Algebraic points and Jónsson's algebraic elements}
\label{sec:p03-w03-p03-033-jonsson-algebraic}

Morley calls the points of \(\operatorname{Tr}^0(A)\) algebraic. He
explains in a footnote that the terms algebraic and transcendental come
from algebraically closed fields of characteristic \(0\), and that his
notion of algebraic bears on a broader concept of algebraic extension
that Jónsson had studied \cite[p.~520, n.~12]{Morley1965}. He
does not say what the relation is. Jónsson works with a class \(K\) closed
under isomorphic copies, subsystems and unions of nonempty directed
families, and having the amalgamation property in the isomorphism form
\cite[pp.~179--180]{Jonsson1962}. By
\cref{thm:p03-w03-p03-033-jonsson-conditions}, \(\mathcal N(T)\) is such
a class. Since its systems carry only relations, the extension \(A(U)\)
of \(A\) by a set \(U\) in Jónsson's sense
\cite[p.~184, Definition~4.1]{Jonsson1962} is the subsystem on
\(|A|\cup U\).

\begin{hbdefinition}{def:p03-w03-p03-033-jonsson-algebraic}{Jónsson's algebraic elements in \(\mathcal N(T)\)}
Let \(A\subseteq B\) and \(A\subseteq C\) in \(\mathcal N(T)\). Following
\cite[p.~185, Definitions~4.3 and 4.4]{Jonsson1962}: an element \(u\) of
\(B\) is \emph{\(A\)-equivalent} to an element \(v\) of \(C\) when some
monomorphism of \(A(u)\) onto \(A(v)\) fixes every element of \(A\) and
maps \(u\) to \(v\); and \(u\) is \emph{algebraic over} \(A\) when every
\(C\supseteq A\) in \(\mathcal N(T)\) has only finitely many elements that
are \(A\)-equivalent to \(u\). The system \(B\) is an \emph{algebraic
extension} of \(A\) when every element of \(B\) is algebraic over \(A\).
\end{hbdefinition}

\begin{hbtheorem}{thm:p03-w03-p03-033-jonsson-algebraic}{proved=annex:p03-w03-p03-033-jonsson-algebraic}{Morley's algebraic points are Jónsson's algebraic elements}
Let \(A\subseteq B\) in \(\mathcal N(T)\), let \(u\in B\), and let
\(p\in S(A)\) be the point realized by \(u\). The following are
equivalent.
\begin{enumerate}
\item \(p\) is algebraic, that is, \(p\in\operatorname{Tr}^0(A)\).
\item \(u\) is algebraic over \(A\) in the sense of
\cref{def:p03-w03-p03-033-jonsson-algebraic}.
\item There is a positive integer \(n\) such that every
\(C\supseteq A\) in \(\mathcal N(T)\) has at most \(n\) elements
\(A\)-equivalent to \(u\).
\item Some formula \(\varphi\in p\) and some \(n\) satisfy
\(T(A)\vdash\neg\exists v_0\cdots\exists v_n
\bigl(\bigwedge_{i\leq n}\varphi(v_i)\wedge\bigwedge_{i<j\leq n}v_i\neq v_j\bigr)\).
\end{enumerate}
Consequently \(B\) is an algebraic extension of \(A\) in Jónsson's sense
exactly when every point of \(S(A)\) realized in \(B\) is algebraic in
Morley's sense.
\end{hbtheorem}

The equivalence of the second and third clauses is Jónsson's Lemma~4.5
\cite[p.~185]{Jonsson1962}, here obtained in \(\mathcal N(T)\) from
compactness. In Morley's Example~I with \(A=\{0\}\), the type of a cube
root of \(2\) contains \(v_0^3=2\), which has three solutions in every
algebraically closed field of characteristic \(0\), so it is algebraic in
both senses; the transcendental point over \(A\) is realized infinitely
often in every model containing \(A\), so it is algebraic in neither sense
(\cref{thm:p03-morley-field-example}). The theorem concerns rank \(0\)
only. Points of positive rank have no counterpart in Jónsson's notion, and
the theorem does not identify the rank of a point with a count of
\(A\)-equivalent elements.

\subsection{Prime models over countable systems}
\label{sec:p03-w03-p03-033-countable-prime}

Morley's prime-model theorem assumes that the isolated points are dense in
\(S(A)\) for every \(A\in\mathcal N(T)\)
(\cref{thm:p03-morley-prime-models}). His footnote to it records that for
countable \(A\) the existence of a model prime over \(A\) had been proved
by Vaught under a somewhat weaker hypothesis
\cite[p.~530, n.~20]{Morley1965}. He does not state the hypothesis. The
next theorem takes a hypothesis on \(A\) alone from Vaught's existence
equivalence (\cref{thm:p03-vaught-isolation-density}) and proves that
Morley's hypothesis implies it. For countable \(A\in\mathcal N(T)\), the
theory \(T(A)\) of \cref{def:p03-morley-type-space} is complete in the
countable language \(L(A)\) of relations and individual constants and has
an infinite model, so it satisfies the standing assumptions of
\cref{def:p03-vaught-standing}; \(F_n(T(A))\) is its algebra of formulas
in \(v_0,\dots,v_{n-1}\), and atomistic is meant in the sense of
\cref{def:p03-vaught-formula-algebra}.

\begin{hbtheorem}{thm:p03-w03-p03-033-countable-prime}{proved=annex:p03-w03-p03-033-countable-prime}{Prime models over a countable system}
Let \(A\in\mathcal N(T)\) be countable.
\begin{enumerate}
\item If \(F_n(T(A))\) is atomistic for every \(n<\omega\), there is a
model \(B\supseteq A\) of \(T\) that is prime over \(A\) in the sense of
\cref{def:p03-morley-prime}, and \(B\) is countable.
\item If the isolated points are dense in \(S(A')\) for every
\(A'\in\mathcal N(T)\), then \(F_n(T(A))\) is atomistic for every
\(n<\omega\).
\end{enumerate}
\end{hbtheorem}

The hypothesis of the first clause concerns only formulas over \(A\);
Morley's hypothesis concerns every system of \(\mathcal N(T)\). The second
clause is the sense in which the first hypothesis is weaker. Whether it is
strictly weaker for some theory is not decided here, and the footnote does
not say.

\subsection{A theory with no countably saturated model in a given power}
\label{sec:p03-w03-p03-033-coding-example}

In the footnote to his Theorem~5.2 Morley asserts that one can give an
example of a theory that is not totally transcendental, and an uncountable
\(\kappa\) with \(\kappa^{\aleph_0}\neq\kappa\), such that the theory has
no countably saturated model of power \(\kappa\); he does not give it
\cite[p.~535, n.~24]{Morley1965}. The first statement of that footnote,
existence when \(\kappa^{\aleph_0}=\kappa\), is
\cref{thm:p03-w03-p03-034-mv-existence}, and
\cref{thm:p03-w03-p03-034-continuum-bound} gives an example with
\(\kappa=\aleph_1\) when \(2^{\aleph_0}>\aleph_1\). The example below
needs no hypothesis beyond the axioms of set theory with Choice. As in
\cref{sec:p03-w03-p03-034-countable-saturation}, countably saturated
means saturated over countable subsystems (\cref{def:p03-morley-saturation}).

\begin{hbexample}{def:p03-w03-p03-033-coding-theory}{A theory that codes sequences}
Fix a bijection \(\langle\cdot,\cdot\rangle:\omega\times\omega\to\omega\).
Let \(Y=\omega^\omega\), the set of functions from \(\omega\) to
\(\omega\), and for \(n<\omega\) let \(f_n:Y\to Y\) be given by
\(f_n(x)(m)=x(\langle n,m\rangle)\). Let \(\Sigma_{\mathrm{cd}}\) be the
complete theory of the structure \((Y,f_n)_{n<\omega}\), in the countable
language with one unary function symbol for each \(f_n\), and let
\(T_{\mathrm{cd}}=\Sigma_{\mathrm{cd}}^*\), in which the graph of \(f_n\)
is a relation \(G_n\) as in \cref{def:p03-morley-standing}.
\end{hbexample}

\begin{hbtheorem}{thm:p03-w03-p03-033-coding-theory}{proved=annex:p03-w03-p03-033-coding-theory}{No countably saturated model of power \(\aleph_\omega\)}
For the theory of \cref{def:p03-w03-p03-033-coding-theory}:
\begin{enumerate}
\item \(T_{\mathrm{cd}}\) is not totally transcendental; indeed \(S(A)\)
has power \(2^{\aleph_0}\) for some two-element
\(A\in\mathcal N(T_{\mathrm{cd}})\).
\item If a model \(B\) of \(T_{\mathrm{cd}}\) of uncountable power
\(\lambda\) is saturated over countable subsystems, then
\(\lambda^{\aleph_0}=\lambda\).
\item The cardinal \(\kappa=\aleph_\omega\) satisfies
\(\kappa^{\aleph_0}\neq\kappa\). The theory has models of power
\(\aleph_\omega\), and none of them is saturated over countable
subsystems.
\end{enumerate}
Consequently, for uncountable \(\kappa\), \(T_{\mathrm{cd}}\) has a model
of power \(\kappa\) saturated over countable subsystems if and only if
\(\kappa^{\aleph_0}=\kappa\).
\end{hbtheorem}

The mechanism is that one element codes a whole \(\omega\)-sequence of
elements through the functions \(f_n\), so a countably saturated model
must contain a separate element for every \(\omega\)-sequence of its
elements. The example uses no ordering and is independent of the
dense-order case of \cref{thm:p03-w03-p03-034-continuum-bound}. It does
not conflict with Morley's Theorem~5.2: that theorem assumes total
transcendence, which fails here by the first clause.
